\label{sec:discussion}

The results support four main conclusions.

\paragraph{Constitutive AD is the preferred derivative route on the flagship
Plasticity3D case.}
\Cref{fig:plasticity3d-derivative-ablation,tab:plasticity3d-derivative-ablation}
show a clean cost-quality picture on the locked
$P_4(L_1)$, $\lambda_{\mathrm{sr}}=1.5$, 8-rank case. All three routes
reach essentially the same terminal state with the same nonlinear and Krylov
counts, so the comparison is about cost rather than about different converged
solutions. On that evidence, constitutive AD is the preferred maintained route:
it preserves the scalar-energy implementation contract while reducing wall time
substantially relative to full element AD and colored SFD.

\paragraph{Validation has to be separated by claim level.}
The Plasticity3D validation ladder in
\Cref{fig:plasticity3d-validation,tab:plasticity3d-validation} changes how the
paper should phrase its mechanics claims. Same-case source-faithfulness to the
source-style surrogate implementation is strong, and that is an important
result for software credibility. The fixed-lambda source-operator comparison
also satisfies the configured endpoint field-agreement thresholds once the
source boundary condition is matched to the maintained glued-bottom mask. The
right interpretation is therefore stronger but still scoped: the maintained
implementation faithfully reproduces the intended endpoint-surrogate workflow,
yet the present evidence does not justify describing that surrogate as a
path-consistent incremental elastoplastic history solve or as a verified
numeric reproduction of the Sysala source-family papers.

\paragraph{Matched external baselines are useful when the contract is narrow.}
The hyperelastic companion comparison against JAX-FEM in
\Cref{fig:hyperelasticity-jaxfem-baseline,tab:hyperelasticity-jaxfem-baseline}
shows that a narrow, fairness-first external baseline can survive review-level
scrutiny. Here the mesh and prescribed displacement schedule are matched, and
the terminal states are post-compared with the repository energy; the resulting
energy and field differences are very small. The comparison is useful precisely
because it is narrow and hyperelastic-only. The paper should not generalize it
into a broad ranking between frameworks or into plasticity validation, but it
does strengthen the manuscript by showing that at least one recent external
implementation can be matched on a shared hyperelastic problem contract.

\paragraph{Fairness and limitations remain part of the scientific result.}
Several limitations should stay explicit. The promoted Plasticity3D strong
scaling study at $\lambda_{\mathrm{sr}}=1.0$ is historical workflow
evidence, not a direct continuation of the glued-bottom
$\lambda_{\mathrm{sr}}=1.55$ matched-comparison figures. The topology study is
reported as end-to-end workflow scaling rather than fixed-work strong scaling.
Some source/local PMG studies remain solver-engineering diagnostics rather than
fully symmetric scientific baselines. None of these caveats erase the paper's
main contribution, but they do define its proper scope: a maintained
software-and-methods workflow with explicit derivative-route choices, explicit
validation surfaces, and explicit fairness boundaries.

The broader methodological lesson is that differentiable FEM becomes most
scientifically useful when local differentiation and sparse solver engineering
are designed together. Automatic differentiation alone does not settle the
behavior of large nonlinear FEM workflows; globalization, sparse ownership,
preconditioning, and the interpretation of surrogate models matter just as
much. Within that scoped claim, \repo{} contributes a reproducible maintained
workflow rather than a universal differentiable-PDE platform or a new
constitutive theory.
